\documentclass[12pt,a4paper]{article}
\usepackage[utf8]{inputenc}
\usepackage[T2A]{fontenc}
\usepackage[russian]{babel}
\usepackage{amsmath, amssymb}
\usepackage{graphicx}
\usepackage{geometry}
\usepackage{listings}
\usepackage{xcolor}
\usepackage{setspace}

\geometry{left=2cm,right=2cm,top=2cm,bottom=2cm}

\lstset{
    language=C++,
    basicstyle=\ttfamily\small,
    keywordstyle=\color{blue},
    commentstyle=\color{green!60!black},
    stringstyle=\color{red},
    numbers=left,
    numberstyle=\tiny\color{gray},
    frame=single,
    breaklines=true,
    showstringspaces=false
}

\begin{document}

\pagestyle{empty}
\begin{center}
    \setstretch{1.1}
    {\fontsize{12pt}{14pt}\selectfont
        МИНИСТЕРСТВО НАУКИ И ВЫСШЕГО ОБРАЗОВАНИЯ\\
        РОССИЙСКОЙ ФЕДЕРАЦИИ\\
        Федеральное государственное бюджетное образовательное учреждение\\
        высшего образования\\
        \textbf{АДЫГЕЙСКИЙ ГОСУДАРСТВЕННЫЙ УНИВЕРСИТЕТ}\\
    }
    \vspace{0.3cm}
    {\fontsize{12pt}{14pt}\selectfont
        Инженерно-физический факультет\\
        Кафедра автоматизированных систем обработки информации и\\
        управления\\
    }
\end{center}
\vspace{4.5cm}
\begin{center}
    \setstretch{1.3}
    {\fontsize{14pt}{16pt}\selectfont ОТЧЕТ ПО ПРАКТИКЕ\\}
    \vspace{0.5cm}
    {\fontsize{16pt}{18pt}\selectfont \textbf{Программная реализация численного метода}\\}
    {\fontsize{14pt}{16pt}\selectfont \textit{Текст из задания по варианту 8}\\}
    \vspace{0.8cm}
    {\fontsize{12pt}{14pt}\selectfont 1 курс, группа УТС\\}
\end{center}
\vspace{3cm}
\hfill
\begin{minipage}{0.48\textwidth}
    \setstretch{1.2}
    {\normalsize Выполнил:\\
    \underline{\hspace{3.5cm}} Д. К. Уваров\\
    «\underline{\hspace{0.8cm}}» \underline{\hspace{2.2cm}} 2026 г.\\
    \vspace{0.5cm}
    Руководитель:\\
    \underline{\hspace{3.5cm}} С. В. Теплоухов\\
    «\underline{\hspace{0.8cm}}» \underline{\hspace{2.2cm}} 2026 г.}
\end{minipage}

\vspace{1.5cm}

\begin{center}
    {\normalsize Майкоп, 2026 г.}
\end{center}

\newpage

\pagestyle{plain}
\setcounter{page}{2}

\section{Цель работы}

Реализовать алгоритм численного вычисления первой производной функции в точке с использованием трёх двухточечных методов.

\section{Теоретические сведения}

Производная функции $y(x)$ в точке $x_0$ равна тангенсу угла наклона касательной к графику функции в этой точке:
\[
f'(x_0) = \tan(\alpha)
\]

Численное дифференцирование — это аппроксимация производной с использованием значений функции в соседних точках.

\subsection{Формулы методов}

Пусть $h$ — малое приращение аргумента. Тогда:

\begin{itemize}
    \item \textbf{Метод 1 (прямая разность, вправо)}:
    \[
    f'(x) \approx \frac{f(x+h) - f(x)}{h}
    \]
    
    \item \textbf{Метод 2 (обратная разность, влево)}:
    \[
    f'(x) \approx \frac{f(x) - f(x-h)}{h}
    \]
    
    \item \textbf{Метод 3 (центральная разность)}:
    \[
    f'(x) \approx \frac{f(x+h) - f(x-h)}{2h}
    \]
\end{itemize}

\subsection{Погрешность методов}

Центральная разность даёт более точный результат, так как погрешность имеет порядок $O(h^2)$, в то время как односторонние методы — $O(h)$.

\section{Реализация}

Программа написана на языке C++. Основные функции:

\begin{lstlisting}
// Прямая разность (вправо)
double derivativeForward(const Function& f, double x, double h) {
    return (f(x + h) - f(x)) / h;
}

// Обратная разность (влево)
double derivativeBackward(const Function& f, double x, double h) {
    return (f(x) - f(x - h)) / h;
}

// Центральная разность
double derivativeCentral(const Function& f, double x, double h) {
    return (f(x + h) - f(x - h)) / (2.0 * h);
}
\end{lstlisting}

\section{Результаты тестирования}

\subsection{Функция $f(x) = x^2$ в точке $x = 2$}

Точное значение производной: $f'(2) = 4$

При $h = 0.001$ получены следующие результаты:

\begin{table}[h!]
\centering
\begin{tabular}{|c|c|c|c|}
\hline
\textbf{Метод} & \textbf{Значение} & \textbf{Абсолютная погрешность} & \textbf{Относительная погрешность} \\
\hline
Прямая разность & 4.001 & 0.001 & 0.025\% \\
Обратная разность & 3.999 & 0.001 & 0.025\% \\
Центральная разность & 4.000 & $3.7 \times 10^{-13}$ & $9.26 \times 10^{-12}\%$ \\
\hline
\end{tabular}
\end{table}

\subsection{Анализ точности}

Как видно из таблицы, метод центральной разности даёт значительно меньшую погрешность по сравнению с односторонними методами.

\section{Выводы}

В ходе выполнения работы:

\begin{enumerate}
    \item Реализованы три метода численного дифференцирования.
    \item Проведено сравнение точности методов на тестовой функции $f(x) = x^2$.
    \item Установлено, что метод центральной разности является наиболее точным.
    \item Программа включает проверку входных данных и обработку ошибок.
    \item Предусмотрен интерактивный режим для работы с различными функциями.
\end{enumerate}

\end{document}